\documentclass[10pt,a4paper]{amsart}
\usepackage[margin=19mm]{geometry}
\usepackage{amsmath,amssymb,mathtools}
\usepackage[scr=rsfs,cal=euler]{mathalfa}
\usepackage{xfrac}
\usepackage[unicode,hidelinks,bookmarks=false]{hyperref}
\newcommand{\ie}{i.e.\ }
\newcommand{\cf}{cf.\ }
\newcommand{\resp}{respectively\ }
\newcommand{\Subsection}{\subsection}
\providecommand{\nobreakdash}{\nobreak\hbox{-}}
\setlength{\emergencystretch}{2em}
\multlinegap0pt
\usepackage{float}
\usepackage{xcolor}
\usepackage{amssymb}
\usepackage{tikz}
\usepackage{enumerate}
\usepackage{bm}
\usepackage{mathtools}
\usepackage{xfrac}
\newtheorem{definition}{Definition}
\newtheorem{lemma}{Lemma}
\newcommand{\sal}{\mathbb{S}_{\Gamma}} %Salvetti complex
\newcommand{\salb}[1]{\mathbb{S}_\Gamma^{#1}} %blowup of a Salvetti - changed to this version for the symmetric spine paper
\newcommand{\spine}{K_{\Gamma}} %spine of untwisted Outer space
\title{Realising VCD for untwisted automorphism groups of RAAGs}
\author{Gabriel Corrigan}
\date{Source: arXiv:2501.03900v2 (CC BY 4.0)}
\begin{document}
\maketitle

\noindent\emph{Source and licence.} Realising VCD for untwisted automorphism groups of RAAGs. Version arXiv:2501.03900v2 (CC BY 4.0), distributed by arXiv under the Creative Commons Attribution 4.0 International licence, http://creativecommons.org/licenses/by/4.0/, as recorded in the arXiv licence metadata for this exact version and shown on the abstract page https://arxiv.org/abs/2501.03900v2. Full licence notice: \texttt{licenses/arxiv-2501.03900-CC-BY-4.0.txt}.\par\smallskip

\noindent\textit{Editorial scope.} Counted unit: the article's lemma (n-1)-face (source0.tex lines 1906--1908). The article's complete original proof of it is delivered with it (source0.tex lines 1910--1925, one \texttt{proof} environment of 16 lines). No mathematics is written, repaired or re-derived: the statement and the proof are the article's own lines, in the article's own notation, and no retained line is substituted or re-flowed.  Retained uncounted as the source-local context the proof needs: lines 1098--1100.  Nothing was omitted from any retained span, so the package carries no omission record.  No retained line contains a citation, so the wrapper carries no bibliography and no bibliography entry.  No retained line contains a figure environment, an image-inclusion command or drawing code, so no image is delivered and the package declares no asset dependency.  Editorial formatting only, in addition: an AMS \texttt{amsart} preamble that restates the article's own declarations and macros used by the retained lines, namely the environments \texttt{definition}, \texttt{lemma} on separate counters, together with the standard packages those lines need.  The retained environments print in this document as Definition 1, Lemma 1 in order of appearance, whereas the article prints the same items with the numbering of its own full document.  The complete native source of the article as distributed in the arXiv e-print for this exact version is bundled unchanged as \texttt{sources/171281/source0.tex}.
\begin{definition}\label{def:untwisted spine}
    The \emph{spine}~$\spine$ is the geometric realisation of the poset of marked~$\Gamma$-complexes.
\end{definition}

\begin{lemma}\label{lem:meet on (n-1)-face}
    Let~$\Gamma$ be a simplicial graph and~$\Pi$ be an inextendible set of~$\Gamma$-partitions. Let~$\alpha$,~$\beta$ be (distinct, non-equivalent) markings of the~$\Gamma$-complex~$\salb{\Pi}$. Suppose that the cubes~$\mathcal{C}_{\alpha} := c(\emptyset, \Pi; \alpha)$ and~$\mathcal{C}_{\beta} := c(\emptyset, \Pi; \beta)$ meet along some common face. Then this face is contained in a face common to both cubes which contains the vertex corresponding to the maximal blowup,~$\salb{\Pi}$.
\end{lemma}

\begin{proof}
Let~$\Phi$ be a subset of~$\Pi$ such that the blow-up along~$\Phi$ is a vertex contained in the common face of~$\mathcal{C}_{\alpha}$ and~$\mathcal{C}_{\beta}$. If~$\Phi = \Pi$, we are done, so assume not.

Let~$c : \salb{\Pi} \to \salb{\Phi}$ be the collapse along the hyperplane set~$\Pi \setminus \Phi$. Let~$\alpha_{\Phi} : \salb{\Phi} \to \sal$ be the collapse to the Salvetti in the cube~$\mathcal{C}_{\alpha}$ and similarly define~$\beta_{\Phi}$ (note that these are necessarily collapses to different marked Salvettis, since~$\alpha$ and~$\beta$ were non-equivalent). 

We have that~$(\salb{\Phi}, \alpha_{\Phi})$ and~$(\salb{\Phi}, \beta_{\Phi})$ are equivalent marked~$\Gamma$-complexes (recall \S\ref{def:untwisted spine}), so there exists an isomorphism~$h_{\Phi} : \salb{\Phi} \to \salb{\Phi}$ such that~$\alpha_{\Phi} \simeq \beta_{\Phi} \circ h_{\Phi}$.

Since~$(\salb{\Phi}, \alpha_{\Phi})$ is obtained from~$(\salb{\Pi}, \alpha)$ by a hyperplane collapse, we have~$\alpha_{\Phi} = \alpha \circ c^{-1}$. Similarly,~$\beta_{\Phi} = \beta \circ c^{-1}$. Hence,
\begin{eqnarray*}
    \alpha & = & \alpha_{\Phi} \circ c \\
    & \simeq & (\beta_{\Phi} \circ h_{\Phi}) \circ c \\
    & = & \beta \circ c^{-1} \circ h_{\Phi} \circ c.
\end{eqnarray*}

So, defining~$h:\salb{\Pi} \to \salb{\Pi}$ as~$h := c^{-1} \circ h_{\Phi} \circ c$, we see that in fact~$(\salb{\Pi}, \alpha) \sim (\salb{\Pi}, \beta)$, via this isomorphism~$h$. Thus the common face of~$\mathcal{C}_{\alpha}$ and~$\mathcal{C}_{\beta}$ must contain this marked~$\Gamma$-complex.
\end{proof}

\end{document}
